% sections/01_introduction.tex
\chapter{Introduction}
\label{ch:introduction}

The generation of organic waste has become one of the defining environmental challenges of our time, particularly as urbanization and industrial activity continue to expand. A major contributor is \glsauto{fw}, a global stream that in the European Union alone is estimated at \qty{88}{\mega\tonne\per\yr}, excluding inedible parts removed from the supply chain \parencite{scherhaufer2018}. When considering only post-processing activities (wholesale, retail, food service and households) approximately \qty{66.3}{\mega\tonne} of \glsauto{fw} are produced annually in the EU-27+1 \parencite{albizzati2021}. The scale of this loss highlights the urgent need for improved management strategies that move beyond simple disposal.

The environmental burden of \glsauto{fw} is substantial. It accounts for roughly \qtyrange{15}{16}{\percent} of the total climate impact of the entire food supply chain, translating to about \qty{186}{\mega\tonne\CO}-eq of greenhouse gas emissions, alongside significant contributions to acidification and eutrophication \parencite{scherhaufer2018}. The majority of these impacts arise during primary production, making prevention the single most effective strategy; however, when waste cannot be avoided, its valorisation through anaerobic digestion or composting is favoured over landfill disposal \parencite{albizzati2021, scherhaufer2018}. Despite this, much \glsauto{fw} continues to be landfilled or incinerated, forfeiting the opportunity to recover its embedded energy and nutrients.

Parallel to the \glsauto{fw} challenge, the rapid expansion of biodiesel production has generated an excess of another organic by-product: \glsauto{cgl}. Formed in the transesterification of triglycerides, approximately \qty{10}{\kilo\gram} of \glsauto{cgl} are produced for every \qty{100}{\kilo\gram} of biodiesel \parencite{garlapati2016}. Global biodiesel output, driven by renewable energy mandates, has caused the \glsauto{gl} market to swell with \glsauto{cgl} supply projected to reach about \qty{6.33}{\mega\tonne} by 2025, of which nearly \qty{680}{\kilo\tonne} are expected to be highly impure due to the shift towards waste-based feedstocks \parencite{attarbachi2023}.

\Glsentrylong{cgl} composition varies widely depending on the oil source and catalyst used, but it typically contains methanol, soaps, free fatty acids, salts and other \glsauto{mong} compounds \parencite{anand2012, attarbachi2023}. These impurities are known to inhibit microbial growth and fermentation, as demonstrated by the severe reduction in \glsauto{pdo} production when untreated \glsauto{cgl} from jatropha, soybean or sunflower oils was used \parencite{anand2012}. Purification to technical or pharmaceutical grade is energy-intensive and often uneconomical for small and medium biodiesel plants \parencite{attarbachi2023}, prompting the search for direct, cost-effective valorisation routes. Because of its high organic content and ready biodegradability, \glsauto{cgl} has attracted attention as an ideal co-substrate for anaerobic digestion, where it can serve as a concentrated carbon source while being simultaneously detoxified by dilution and microbial synergy \parencite{gebreegziabher2025, tomczak2025}.

\Glsentrylongauto{ad} is a mature biological process in which consortia of microorganisms convert organic matter into biogas, primarily \glsauto{ch4} (\qtyrange{55}{65}{\percent}) and \glsauto{co2}, under \glsauto{o2}-free conditions, passing through the sequential stages of hydrolysis, acidogenesis, acetogenesis and methanogenesis \parencite{gebreegziabher2025}. \Glsentrylong{ad} offers a unique combination of benefits: it produces renewable, storable energy that can replace fossil fuels; generates a nutrient-rich digestate that can be used as organic fertiliser; and reduces the environmental load of organic wastes by avoiding landfilling and its associated \glsauto{ch4} emissions \parencite{scherhaufer2018, gebreegziabher2025}.

Within this framework, anaerobic co-digestion, the simultaneous treatment of two or more substrates, has emerged as a particularly effective strategy. Co-digesting \glsauto{cgl} with nitrogen-rich materials such as animal manure, sewage sludge or \glsauto{fw} balances the \glsauto{cn_ratio}, dilutes potential inhibitors, provides essential macro and micro nutrients, and fosters synergistic microbial interactions \parencite{tomczak2025, gebreegziabher2025}. The result can be a dramatic improvement in biogas and \glsauto{ch4} yield: studies report increases ranging from \qtyrange{14}{226}{\percent} when \glsauto{cgl} is added to a variety of base substrates, with \glsauto{ch4} contents generally exceeding \qty{60}{\percent} as long as \glsauto{gl} does not exceed about \qty{8}{\percent} (\glsauto{vpv}) of the feed \parencite{gebreegziabher2025, tomczak2025}.

However, excessive \glsauto{gl} loads can lead to \glsauto{vfa} accumulation, pH drop and methanogen inhibition, underscoring the need for careful optimisation \parencite{gebreegziabher2025, tomczak2025}. Overall, the convergence of abundant \glsauto{fw} and surplus \glsauto{cgl} supplies, together with the proven capability of \glsauto{ad} to convert these streams into renewable \glsauto{ch4} and fertiliser, represents a tangible route towards a circular, low-carbon bioeconomy.

\chapter{Objectives}
\label{ch:objectives}

\section{General Objectives}
\label{sec:general_objectives}

The primary objective of this study was to evaluate the efficiency of a single-stage anaerobic co-digestion system using organic \glsauto{fw} and \glsauto{cgl} as substrates. The investigation will be conducted by determining biogas production, with a focus on industrial feasibility.

The central hypothesis is that the controlled addition of \glsauto{cgl}, a by-product of the biodiesel industry, helps biogas production due to its high energy content provided that adequate operational conditions are maintained to avoid \glsauto{vfa} overload.

\section{Specific Objectives}
\label{sec:specific_objectives}

\begin{enumerate}
    \item Characterize the substrates (\glsauto{fw} and \glsauto{cgl}), evaluating physicochemical parameters relevant to anaerobic digestion.
    \item Operate the single-stage anaerobic co-digestion reactor, assessing process stability.
    \item Efficiency evaluation of the process for the utilization of different substrates and different \glsauto{olr}.
\end{enumerate}

